\begin{korollar}{Maximum zentrierter normalverteilter Zufallsgrößen}
             {MaximumZentrierterNormalverteilterZufallsgroessen}
Seien $n \in \N$, $\sigma \in \left(\R_{>0}\right)^n$,
$\gamma: \Omega \rightarrow \R^n$, $\delta: \Omega \rightarrow \R^n$ mit
$\delta$ unabhängig und
\[\forall\, i \in \haken n: \gamma_i, \delta_i \text{ sind }
N(0,\sigma_i^2)\text{--verteilt.}\]
Dann gilt:
\[\E\left(\max_{i \in \haken n}\gamma_i \right)
\le \sqrt 2 \E\left(\max_{i \in \haken n}\delta_i\right)\]
und diese Ungleichung ist optimal.
\end{korollar}
\begin{proof}
Seien $i, j \in \haken n$.\\
Dann gilt $\sigma_i^2 - 2\sigma_i\sigma_j +\sigma_j^2
= (\sigma_i - \sigma_j)^2 \ge 0$ und damit $2 \sigma_i\sigma_j \le
\sigma_i^2 + \sigma_j^2$, also
\begin{align}
(\sigma_i + \sigma_j)^2 = \sigma_i^2 + 2 \sigma_i\sigma_j + \sigma_j^2
\le \sigma_i^2 + \sigma_i^2 + \sigma_j^2 + \sigma_j^2 
= 2(\sigma_i^2 + \sigma_j^2) \text. \label{mzngeq:1}
\end{align}
Außerdem gilt:
\begin{align*}
\norm{\gamma_i - \gamma_j}_2^2 
&\le \left(\norm{\gamma_i}_2 + \norm{\gamma_j}_2 \right)^2
= (\sigma_i + \sigma_j)^2 \overset{\eqref{mzngeq:1}}\le 
2(\sigma_i^2 + \sigma_j^2)\\
&= 2\left(\norm{\delta_i}_2^2 + \norm{\delta_j}_2^2 \right)
\overset{\delta \text{ unabh.}}=
2 \left(\norm{\delta_i}_2^2 + \norm{\delta_j}_2^2 
- 2 \sprod{\delta_i}{\delta_j}_{L_2} \right)\\
&= 2\norm{\delta_i - \delta_j}_2^2 =
\norm{\sqrt 2 \delta_i - \sqrt 2 \delta_j}_2^2\text.
\end{align*}
Offenbar ist für jedes $k\in \N$ die Abbildung $\sqrt 2\delta_k$
$N(0, 2\sigma_k^2)$--verteilt. Damit sind die Voraussetzungen von
\ref{SatzVonFernique} für $\gamma$ und $\sqrt 2 \delta$ erfüllt, sodass gilt:
\[\E\left(\max_{i \in \haken n}\gamma_i \right)
\overset{\ref{SatzVonFernique}}\le 
\E\left(\max_{i \in \haken n}\sqrt 2\delta_i\right)
\overset{\sqrt 2>0}=\sqrt 2\E\left(\max_{i \in \haken n}\delta_i\right)\text.\]
Für den Beweis der Optimalität seien $Y_1, Y_2$ stochastisch unabhängige
$N(0,1)$--verteilte Zufallsgrößen, sowie $X_1 :=Y_1$ und $X_2:=-Y_1$.\\
Dann sind auch $X_1, X_2$ $N(0,1)$--verteilt, aber nicht stochastisch
unabhängig.\\
Offensichtlich ist $\max\meng{X_1,X_2} = \max\meng{Y_1, -Y_1} = \abs {Y_1}$
und damit
\begin{align}
\E\left(\max_{i \in \haken 2}X_i \right) = \E\left(\abs {Y_1} \right) 
= \sqrt{\frac 2 \pi}\label{mzngeq:2}\text.
\end{align}
Nach \cite{K2} (dort Abschnitt 6.6) gilt:
\begin{align}
\int_{-\infty}^\infty \exp(-t^2)dt = 2 \int_0^\infty \exp(-t^2)dt = \sqrt \pi
\label{mzngeq:3}\text.
\end{align}
Seien $\Phi$ die Verteilung und $\vi$ die Dichte von $N(0,1)$, sowie
$f$ die Dichte von $\max\meng{Y_1, Y_2}$.\\
Nach \refclaim{EigenschaftenVonOrdnungsstatistiken}
{EigenschaftenVonOrdnungsstatistiken4} gilt dann: $f = 2\Phi\vi$,
sodass wir wegen $\forall\, t \in \R: \vi'(t) = -t\vi(t)$ erhalten:
\begin{align}
\E\left(\max_{i \in \haken 2}Y_i \right)&=
\int_{-\infty}^\infty t f(t)dt = 2\int_{-\infty}^\infty \Phi(t)t\vi(t)dt
= - 2\int_{-\infty}^\infty \Phi(t)\vi'(t)dt\notag\\
&= -2\left([\Phi(t)\vi(t)]_{t=-\infty}^\infty - 
\int_{-\infty}^\infty \vi(t)^2dt \right)
= 2 \int_{-\infty}^\infty \frac 1 {2\pi}\exp\left(-t^2\right)dt\notag\\
&= \frac 1 \pi \int_{-\infty}^\infty \exp\left(-t^2\right)dt
\overset{\eqref{mzngeq:3}}= \frac 1 \pi \sqrt \pi = \sqrt{\frac 1 \pi}\text.
\label{mzngeq:4}
\end{align}
Also
\[\E\left(\max_{i \in \haken 2}X_i \right) \overset{\eqref{mzngeq:2}}= 
\sqrt 2 \sqrt{\frac 1 \pi} \overset{\eqref{mzngeq:4}}= 
\sqrt 2 \E\left(\max_{i \in \haken 2}Y_i \right)\text.\]
\end{proof}